% P10-Core-Inhaltsprüfung, 15.09.2026: datierter Ergebnisanschluss; B-60 korrigiert.
% Belege: thesis-evidence/w2/core-content-audit-20260915/evaluation/
% B-59, 15.09.2026: datierte Ledger-Nachauswertung, Stützung und Reichweite.
% Beleg: thesis-evidence/w2/ledger-repetitions-20260915/ (Originale unverändert).
% P10-HITL, 15.09.2026: Entscheidungszweck und belegte Zustandswirkung präzisiert.
% P10-4c, 14.09.2026: Ledger-Herleitung und Nachweisanschlüsse präzisiert.
% P10-4b, 14.09.2026: Begriffs-/Beleganschlüsse und abgegrenzte Prüfaussagen präzisiert.
% P10-4a, 14.09.2026: Mechanismus, Quellenreichweite und v4-Nachweis angeglichen.
% M19 (12.09.2026): Leserführung und präzise Ergebnisdarstellung;
% Messwerte und historische Nachweisgrenzen bleiben erhalten.
% ============================================================
% Kapitel 8 / 8.1 — kompakte Arbeitsfassung, Phase 02b, 10.09.2026.
% Kapitelkommando wie in 5.1/6.1 in der ersten Abschnittsdatei.
% Grundlage: Kapitel 3–7, M02-Stand; Schlussgliederung-Kapitel-8.md.
% Literatur: Spijkman et al. 2022, S. 281/283; iReDev v1, S. 3,
% gegen lokale Originaltexte geprüft. Bestehende Bib-Schlüssel aus 2.5.
% Die interpretierende Ledger-Passage aus 7.7 ist hier verdichtet;
% Ergebnisprofil, Aufwandsmessung und lokale Grenzen bleiben dort.
% ============================================================

\chapter{Diskussion und Fazit}
\label{chap:diskussion-fazit}

\section{Diskussion der Ergebnisse und des Beitrags}
\label{sec:diskussion-beitrag}

TF1 richtet sich auf die aus Nutzungsziel und Exploration gewonnenen
Designanforderungen (Abschnitt~\ref{sec:einleitung-forschungsfragen}).
Die Entwicklung konkretisierte insbesondere das Problem, aus
wiederkehrender Projektkommunikation fachliche Inhalte zu gewinnen,
deren Herkunft und spätere Änderungen prüfbar bleiben.
Die in Kapitel~\ref{chap:explorative-problemklaerung} rekonstruierte
Exploration zeigte Inhaltslücken in erzeugten Artefakten und
Schwierigkeiten, diese durch einen nachgelagerten Review verlässlich
festzustellen. Der Übergang zur vorgelagerten Quellenstrukturierung
setzte deshalb an der Erzeugung an: Quellenbezüge und Zwischenstände
sollten bereits während der Verarbeitung entstehen. Die frühen
Prüfvergleiche und Evidence-first-Vorstudien begründeten diese
Designrichtung unter ihren jeweiligen Bedingungen. Die damaligen
Modell- und Prototypstände gehören zu dieser Entstehungssituation;
ob ein anderer Gesamtaufbau mit stärkeren Modellen dieselben
Eingriffe entbehrlich gemacht hätte, wurde nicht geklärt.

Die Antwort auf TF2 liegt in der daraus entwickelten Aufgabenverteilung.
Sie begrenzt Autonomie anhand des jeweiligen Entscheidungsgegenstands.
Modelle interpretieren Aussagen,
erzeugen fachliche Vorschläge und treffen bestandsbezogene Einordnungen;
ausgewählte Agenten wählen dafür Werkzeuge. Festgelegte Übergänge und
formal prüfbare Integritätsregeln werden durch Workflows und deterministische
Komponenten umgesetzt. Die fachliche Übernahme in den Fortschreibungspfaden
verbleibt bei menschlichen Entscheidungen. Ein Executor bildet dabei die
technische Einbindung einer Verarbeitungseinheit, deren Entscheidungsspielraum
gesondert festgelegt wird
(Abschnitt~\ref{sec:realisierung-agent-executor-workflow}). Der erfolgreiche
Abschluss einer Stufe richtet sich nach ihren festgelegten
Prüfkriterien; ein beendeter Agentenaufruf genügt dafür nicht
(Abschnitt~\ref{sec:saeule-semantische-transformation}). Die Begründung
für diese Kombination liegt in prüfbaren Übergaben und begrenzten
Zustandswirkungen; eine allgemeine Überlegenheit gegenüber stärker
autonomer Koordination wurde nicht untersucht.

Für TF3 ergibt sich ein gemischtes Ergebnisbild. Der Inhaltsvergleich
zeigt, dass die zusätzliche Strukturierung keinen Vollständigkeitsvorteil
sicherstellt. Im Stressfall erhielt die Ledger-Kette
54 von 109 Referenzaussagen vollständig, die Einzel-Agent-Konfiguration F
dagegen 69 bei geringerem Tokenaufwand
(Abschnitte~\ref{sec:eval-kernzahlen} und~\ref{sec:eval-aufwand}). Auch F
lieferte gültige Quellenreferenzen. Die zusätzliche Verarbeitung geht in
diesen Läufen somit nicht mit einer besseren direkten Abdeckung einher.
Zwei weitere Ledger-Ausgaben ändern diese Einordnung nicht: Ihre
Abdeckung liegt in ähnlicher Größenordnung; welche Inhalte vollständig
erhalten sind, unterscheidet sich jedoch.
Zugleich bildet der Vergleich die
untersuchten Ausgaben ab; eine alternative Realisierung der gesamten
Core-Fortschreibung wurde damit nicht bewertet. Die fehlenden Vertragsfelder
für die bestehenden Folgestufen beschreiben eine konkrete Anschlusslücke,
keine grundsätzliche Unfähigkeit der Vergleichsbedingung
(Abschnitt~\ref{sec:eval-anschluss}).

Eine belegte Leistung des Ledgers liegt in der Untersuchung der
Verarbeitung selbst. Die Quellenverknüpfungen, begründeten
Einordnungen ungenutzter Units und gespeicherten Zwischenstände
adressieren das ursprüngliche Rekonstruktionsproblem: Es lässt
sich nachlesen, welche Bezüge und Prüfentscheidungen die
Verarbeitung tatsächlich festgehalten hat. Die Zwischenstände erlaubten es, zwei
Abdeckungsverluste bei der Kanonisierung zu lokalisieren. Der zugehörige
Übergangs-Check hatte Referenzen geprüft und diese Inhaltsverluste nicht
verhindert (Abschnitt~\ref{sec:eval-stufen}). Die nachträgliche Zuordnung erzeugter Prüfhinweise
zu einem Teil der Abdeckungslücken erschließt außerdem mögliche Prüfansätze;
ihre tatsächliche menschliche Bearbeitung wurde nicht untersucht
(Abschnitt~\ref{sec:eval-luecken}). Einzelne im Aussagentext fehlende Inhalte bleiben in gespeicherten
Belegzitaten erhalten (Anhang~\ref{anh:evaluation},
\hyperref[anh:coverage-nachreview]{Abschnitt~D.8}).
Dies erweitert den prüfbaren Artefaktumfang, erhöht aber nicht die
ausgewiesene Claim-Abdeckung. Nachvollziehbarkeit, Inhaltserhaltung
und erfolgreiche Fehlerbehebung sind somit getrennt zu beurteilen.
Strukturierte Zwischenstände machen bestimmte Fehler untersuchbar;
ihre Erzeugung kann jedoch Inhalte verlieren.

Für den laufübergreifenden Projektbestand zeigen die Fortschreibungsfälle,
dass neue Informationen gegen vorhandene Inhalte eingeordnet und über
Entscheidungen, Relationen und Projektionen weiterverarbeitet werden können
(Abschnitt~\ref{sec:eval-fortschreibung}). Am Besuchsbeispiel wird
dieser Beitrag konkret: Eine spätere GitHub-Eingabe ergänzt eine aus
dem Meeting entstandene Arbeitsplanung und behält ihren eigenen
Eingangsbezug (Abbildung~\ref{fig:core-netz-besuchsbeispiel}).
Agenten übernehmen dabei Interpretation und Vorschlagsbildung;
die festgelegten Anwendungsschritte setzen die freigegebenen
Verbindungen um. Flexible Eingänge und agentische Bedienung nutzen
damit denselben fachlichen Bestand und die vorgesehenen Kontrollen.
Die ergänzende Inhaltsprüfung verfolgt auch einen kenntlich aus
Bestandswissen abgeleiteten Vorschlag bis zur gespeicherten Anforderung,
zum erweiterten Arbeitspaket und zur Issue-Aktualisierung
(Abschnitt~\ref{sec:eval-core-inhalte}). An diesem Fall stimmen
die verbundenen Inhalte fachlich überein; die bloße Existenz ihrer
Referenzen wäre dafür kein hinreichender Nachweis.

Die menschliche Entscheidungshoheit wird dabei an Zustandswirkungen
prüfbar: In F1 entspricht die angewandte Menge der freigegebenen
Auswahl; F3 zeigt die Übernahme einer vom Menschen festgelegten
Konfliktlösung. Die zusätzliche Ablehnung einer Vorlage mit fehlenden
Besuchsstatuswerten zeigt einen weiteren konkreten Eingriff gegen
unvollständige Inhalte. Die Kontrolltests ergänzen diese Fallbelege um
fehlende und abgelehnte Freigaben sowie einen gültigen Übernahmefall
(Abschnitt~\ref{sec:eval-kontrollen}). Damit ist für die untersuchten
Pfade belegt, wie agentische Vorschläge unter menschlicher Kontrolle
wirksam werden. Ob diese Prüfmöglichkeiten die Fehlerhäufigkeit
senken, wurde nicht vergleichend untersucht.

Das verfeinerte Herkunftsaudit belegt im untersuchten Snapshot alle
104 gespeicherten Ledger-Kanten im Ursprungskontext. Für 234 von
237 Items ist mindestens ein direkter Bezug oder gerichteter
Beitragspfad erschlossen; letzterer kann nur Teile des Inhalts
erklären. Gleichzeitig bleiben 64 aktuelle Einzelreferenzen bei
33 Items offen, und nicht jeder jüngste Versionsübergang ist über
die unterstützten Laufbezüge belegt
(Abschnitt~\ref{sec:eval-provenienz}). Der Nachweis stützt somit
konkrete Quellenzugänge und Übernahmeketten, keine lückenlose
Herkunft aller aktuellen Inhalte. Referenzen müssen mit ihrem
Laufkontext verbunden und ihre Zwischenartefakte verfügbar sein.
Die Einordnung einzelner Zitate zeigt zusätzlich, dass ein
auffindbarer Beleg nicht bereits dessen semantische Stützung
bestätigt. Auch die Kontrolltests sichern ausgewählte
Zustandswirkungen, nicht die fachliche Richtigkeit aller Beziehungen.

Diese Befunde präzisieren den Beitrag gegenüber dem Forschungsstand.
TRACE2CONV unterstützt bereits die Rückführung von User Stories auf
relevante Gesprächsbeiträge \cite[S.~281, 283]{spijkman2022back}.
Der in Abschnitt~\ref{sec:forschungsstand-positionierung} ausgewertete
iReDev-Preprint beschreibt spezialisierte Agenten, einen gemeinsamen
Artefaktraum, ereignisgesteuerte Aktualisierungen und menschliche
Eingriffspunkte \cite[S.~3]{jin2025iredev}.
Die vorliegende Arbeit beansprucht entsprechend einen Beitrag durch die
konkrete Verbindung von quellengebundener Verarbeitung, fachlichem
Projektzustand und kontrollierter Fortschreibung sowie deren Untersuchung.
MAF stellt hierfür die in Kapitel~\ref{chap:maf-realisierung} erläuterten
Mechanismen zur Einbindung von Agenten, Werkzeugen und persistierbaren
Workflows bereit. Evidenzmodell, fachliche Identitäten, Relationen und
Governance-Verträge sind projektseitige Gestaltung. Die Realisierung
dokumentiert diese Verbindung mit MAF; einen isolierten Nutzen gegenüber
anderen Frameworks untersucht sie nicht.
